\documentclass[sigconf,nonacm]{acmart}

\newcommand{\NOx}{NO$_x$}
\newcommand{\modelname}{PlumeWatch}

\title{Detecting Power Plant \texorpdfstring{\NOx}{NOx} Emissions Changes using
Physics-Guided Deep Learning and TEMPO Satellite Imagery}

\author{Pranav Walimbe}
\affiliation{%
  \institution{University of California, Berkeley}
  \city{Berkeley}
  \state{California}
  \country{USA}
}
\email{pranav.walimbe@berkeley.edu}

\author{Ron Cohen}
\affiliation{%
  \institution{University of California, Berkeley}
  \city{Berkeley}
  \state{California}
  \country{USA}
}
\email{rcohen@berkeley.edu}

\date{October 3, 2026}

\begin{document}

\begin{abstract}
NASA's Tropospheric Emissions: Monitoring of Pollution (TEMPO) instrument
provides hourly daytime observations of tropospheric nitrogen dioxide
(NO$_2$) across North America. These observations could support independent,
facility-level monitoring of power plant nitrogen oxides (\NOx{}) emissions,
filling gaps caused by continuous emissions monitoring system (CEMS) failures
and covering facilities not subject to hourly emissions-reporting requirements.
Previous studies have estimated hourly point-source emissions from individual
satellite scenes using empirical methods such as cross-sectional flux and the
directional derivative approach. However, these estimates are sensitive to
retrieval uncertainty, changing atmospheric transport, TEMPO's coarse spatial
resolution, signals from nearby emitters, and temporal mismatch with hourly
CEMS totals. We instead formulate facility monitoring as change detection over
sequences of TEMPO observations. This relative target uses repeated scans to
capture plume evolution while reducing sensitivity to persistent background
levels and biases in absolute column values. We develop a physics-guided
convolutional recurrent neural network (CRNN) that processes five consecutive
TEMPO observations, High-Resolution Rapid Refresh (HRRR) meteorological fields,
and transport residuals. The network jointly classifies the target power plant's
emissions as increasing, steady, or decreasing and predicts a spatial map of
contributions from emitting facilities within the observed scene. On a held-out
evaluation set containing more than 20,000 samples from over 100 U.S. power
plants observed between August 2023 and July 2026, the model achieves a
classification accuracy of $X\%$, outperforming a seasonal-only baseline
($Y\%$) and a cross-sectional flux baseline ($Z\%$). These results support the
use of sequential TEMPO imagery for detecting short-term changes in power plant
emissions across facilities and observing conditions.
\end{abstract}

\maketitle

\section{Introduction}

Based on U.S. Environmental Protection Agency (EPA) inventory data, power plants
in the U.S. electricity-generation sector emitted approximately 0.74 million
short tons of \NOx{} in 2023, about 9.5\% of national
emissions~\cite{epa2026emissionsTrends}. \NOx{} emissions contribute to
ground-level ozone and secondary particulate matter, pollutants that can
inflame and damage the airways, reduce lung function, and aggravate
asthma~\cite{epa2016noxHealth}. EPA regulations require most large U.S.
fossil-fuel-fired generating units that produce electricity for sale to submit
unit-level hourly emissions data on a quarterly basis~\cite{epa2026powerMonitoring}.
However, monitoring equipment malfunctions and maintenance can make CEMS
measurements unavailable or invalid, requiring plants to report calculated
substitute values instead of direct measurements~\cite{epa2021substitute}.
Behind-the-meter generators that do not sell electricity can fall outside
federal hourly reporting requirements, as illustrated by the natural-gas
turbines powering SpaceXAI's Colossus 2 data center~\cite{gauld2026dataCenter}.
Outside the United States, comparable public hourly, unit-level emissions data
are generally unavailable, even where facilities continuously monitor
emissions~\cite{huang2026detectability}. Independent estimates could therefore
fill reporting gaps and reveal short-term changes before the EPA's quarterly release
of hourly emissions data.

TEMPO is an ultraviolet-visible spectrometer hosted aboard a geostationary satellite
that observes tropospheric vertical column densities (TVCDs) for various trace gases
across North America during daylight hours~\cite{zoogman2017tempo}. Earlier polar-orbiting
instruments, including the Ozone Monitoring Instrument (OMI) and TROPOspheric Monitoring
Instrument (TROPOMI), typically observe a given location once per day; TEMPO's geostationary
view provides approximately hourly scans with finer spatial
resolution~\cite{levelt2006omi,veefkind2012tropomi,zoogman2017tempo}. Hourly sampling
captures intraday changes in power plant plumes and atmospheric transport that daily
observations can miss. However, power plant signals in TEMPO NO$_2$ TVCDs can be weakened
or obscured by factors including evolving meteorology and atmospheric chemistry, cloud cover,
viewing geometry and retrieval uncertainty, dilution within TEMPO's 2.0-by-4.75 km spatial resolution
per pixel, and interference from nearby emission sources~\cite{tempo2026userGuide,huang2026detectability,gauld2026dataCenter}.
Cross-comparisons with CEMS data, which is the most reliable ground truth available, introduce an
additional temporal mismatch because TEMPO observes a facility at a specific time during each hourly scan,
whereas CEMS data aggregate emissions over an entire clock hour~\cite{tempo2026userGuide,epa2026powerMonitoring}.

Satellite-based point-source studies typically estimate emissions by relating an
observed column enhancement to atmospheric transport. The cross-sectional flux
(CSF) method integrates background-subtracted columns across plume transects, while
the integrated mass enhancement (IME) method converts the total excess plume mass
to a source rate using an estimated residence time~\cite{meier2024conversion,santaren2025benchmarking}.
Gaussian plume methods fit a parameterized dispersion model, whereas flux divergence
and directional derivative approaches (DDA) apply mass-conservation relationships
to satellite columns and meteorological fields~\cite{santaren2025benchmarking,sun2025hourly}.
Santaren et al. compared five inversion methods on synthetic daily satellite images
of three power plants, for which the source emissions were known, and tested the
effects of cloud gaps and wind uncertainty. Gaussian plume and CSF estimates had
the lowest annual-emission errors, while missing observations reduced the number
of usable estimates for all methods~\cite{santaren2025benchmarking}. Sun et al.
applied DDA to individual hourly TEMPO NO$_2$ scans, integrated the resulting fields
over fixed footprints around 14 U.S. power plants, and compared hourly, daily, and
monthly estimates with CEMS. Adding the observed column-tendency term improved
correlations at 10 plants, but the inferred rates remained about four times lower
than CEMS because of retrieval and near-source \NOx{}-to-NO$_2$ conversion
uncertainties~\cite{sun2025hourly}. Beaudry et al. instead fit a rise--decay model
to TEMPO NO$_2$ plumes at 64 plants and compared the resulting average emissions
with CEMS; estimates fell within 50\% of CEMS at 27 plants but showed substantial
low bias below 250~kg NO$_2$~h$^{-1}$~\cite{beaudry2026quantification}.

Machine-learning approaches offer an alternative to prescribing a fixed relationship
between plume structure, atmospheric transport, and source emissions. Prior studies
have addressed different targets. Finch et al. trained a convolutional neural network
on 6,000 labeled TROPOMI NO$_2$ images to classify whether a scene contained a plume,
achieving greater than 90\% test accuracy without estimating its source rate~\cite{finch2022automated}.
Huang and Wang first labeled TROPOMI observations by whether a power plant plume was
detectable, then predicted those labels from plant, meteorological, surface, and sensor
features; their classifiers achieved F1 scores above 0.66 but excluded plants near
major interfering sources~\cite{huang2026detectability}. Alnaim et al. trained recurrent
models on one year of TROPOMI, meteorological, and EPA data to predict plant emissions,
but focused on isolated facilities and did not evaluate generalization to previously
unseen plants~\cite{alnaim2022evaluating}. For other trace gases, U-Plume used a U-Net
plume mask, satellite methane imagery, and wind speed to predict source rates, but its
CNN estimates regressed toward the training-set mean~\cite{bruno2024uplume}. A separate
image-to-scalar CNN trained on simulated CO$_2$ and wind fields reached about 20\%
relative error when evaluated on three held-out power plants, although its evaluation
did not include the retrieval errors present in real scenes~\cite{dumont2024deepLearning}.
To our knowledge, no published study has trained an ML model to infer facility-level
emissions or behavior directly from sequences of TEMPO NO$_2$ imagery.

Estimating absolute hourly emissions from individual TEMPO scenes requires
separating weak facility signals from background pollution and resolving
uncertainties in retrieval, atmospheric chemistry, and transport. We instead
use TEMPO's hourly sampling to detect relative changes over short observation
sequences, reducing dependence on absolute column values and persistent
background levels.

Our contributions are as follows:

\begin{itemize}
    \item We formulate satellite-based power plant monitoring as a three-class
    change-detection task that identifies increasing, steady, and decreasing
    emissions behavior.

    \item We introduce a physics-guided deep learning framework that combines
    sequences of TEMPO NO$_2$ observations with meteorology and transport
    information. The model jointly predicts emissions behavior and a spatial
    map of emissions contributions within the observed scene.

    \item We evaluate the model on held-out power plants across geographic
    regions, seasons, and operating profiles and compare its performance with
    seasonal-only and empirical baselines.
\end{itemize}

\section{Data Collection and Dataset Construction}

\subsection{TEMPO NO\texorpdfstring{$_2$}{2} Data}

NASA distributes TEMPO NO$_2$ retrievals in two formats: Level 2 data on the
instrument's native measurement footprints and Level 3 data regridded from each
east-to-west scan~\cite{tempo2026userGuide}. We downloaded rasters from the version 4 level 2 product
over the months covering August 2023 to August 2026. We chose the Level 2 product
because it records an observation time for each mirror step, allowing us to
later interpolate hourly CEMS records to the approximate time TEMPO scanned each facility.

TEMPO Level 2 observations lie on irregular native footprints, whereas our
model requires a fixed spatial grid. For each power plant, we define a
72~km $\times$ 72~km area of interest (AOI) centered on the facility and divide it into a
$24 \times 24$ grid of 3~km $\times$ 3~km cells. We assign each grid cell
an NO$_2$ TVCD by averaging the valid TEMPO observations that overlap it,
weighted by their area of overlap.

\begin{figure}[t]
    \centering
    \fbox{\parbox[c][1.15in][c]{0.91\columnwidth}{\centering
    Regridding schematic placeholder}}
    \caption{Planned illustration of native TEMPO Level 2 footprints mapped to
    the plant-centered $24 \times 24$ analysis grid.}
    \Description{Placeholder for a schematic showing native TEMPO footprints
    intersecting a plant-centered analysis grid.}
    \label{fig:tempoRegridding}
\end{figure}

During regridding, only native TEMPO pixels with a main data quality flag of
zero, an effective cloud fraction no greater than 0.20, and a finite NO$_2$
value contribute to each grid-cell average. We then create a quality mask
indicating which cells in the regridded raster contain valid NO$_2$ values.
A pretrained model later fills the masked cells.

\subsection{CEMS Data}

We obtained unit-level hourly emissions and operating records and annual
facility and unit attributes from the U.S. Environmental Protection Agency's
Clean Air Markets Program Data (CAMPD) APIs for August 2023 through September
2026.

For each AOI-hour, we summed the unit-level \NOx{} mass reported by all
facilities in the AOI and discarded the hour if any contributing measurement
was invalid or unavailable. We scored each AOI using its median total \NOx{}
mass across hours whose mean unit operating time met or exceeded the
AOI-specific median. We retained the highest-scoring 50\% to focus the dataset
on emissions signals more likely to be detected by TEMPO.

\subsection{HRRR Weather Data}

The High-Resolution Rapid Refresh (HRRR) is an hourly weather model with a
3~km grid covering the contiguous United States. From NOAA's public HRRR
archive, we downloaded hourly 2~m air temperature and eastward and northward
wind components at 80~m above ground for August 2023 through August 2026. We
bilinearly interpolated these fields onto the same $24 \times 24$ pixel AOIs
we compute from CEMS data and TEMPO observations.

\subsection{Sample Construction}

Each sample contains five consecutive observations of one AOI, indexed from
$t_0$ to $t_4$. We retain a sequence only when each pair of adjacent TEMPO
scans is separated by 40--70 minutes and every NO$_2$ raster has valid values
in at least 85\% of its cells. At each timestep, we pair the TEMPO raster with
the HRRR temperature and wind rasters from the closest hour.

Because CEMS records correspond to clock hours rather than TEMPO scan times,
we linearly interpolate the AOI-hour \NOx{} emissions between the surrounding
hourly CEMS records at each scan time. Let $E_i$ denote the interpolated emissions
at time $t_i$. For scans $t_a$ through $t_b$, we calculate the time-weighted mean
\begin{equation}
    \overline{E}_{a:b} =
    \frac{\sum_{i=a}^{b-1}(t_{i+1}-t_i)(E_i+E_{i+1})/2}
    {t_b-t_a}.
\end{equation}
We define the sequence change as
$\Delta E=\overline{E}_{2:4}-\overline{E}_{0:2}$. A sequence is labeled as an
increase when $\Delta E \geq 250~\mathrm{lb\,h^{-1}}$, a decrease when
$\Delta E \leq -250~\mathrm{lb\,h^{-1}}$, and steady otherwise.

To produce transport residuals for each pair of adjacent NO$_2$ rasters, we use a semi-Lagrangian transport model
to advect the earlier raster with the average of the two endpoint wind fields while accounting for
horizontal dispersion and effective NO$_2$ decay. We subtract the transported
estimate from the next TEMPO observation to isolate new emission signal from the later timestep not
captured by the existing plume. We retain the original NO$_2$ quality masks and combine them after
transport to identify unreliable residual pixels.

For each sample, we also compute increase and decrease pixel contribution rasters
(sized $24 \times 24$) to allow the model to learn plume attribution. We first sum unit emissions within
each facility, interpolate each facility's emissions at the five scan times,
and calculate its time-weighted change using the same early and late windows from above.
For an increase sample, positive facility changes are normalized to sum to one
and distributed around their facility locations on the increase map. Decrease
samples use the normalized magnitudes of negative facility changes on the
decrease map. The inactive map is zero, and both maps are zero for steady
samples.

\begin{figure}[t]
  \centering
  \IfFileExists{figures/sample-construction.pdf}{%
      \includegraphics[width=\columnwidth]{figures/sample-construction.pdf}%
  }{%
      \fbox{\parbox[c][1.4in][c]{0.9\columnwidth}{\centering
      \textit{Representative five-timestep sample}}}
  }
  \caption{Representative sample with five NO$_2$ and meteorological
  observations, four transport residuals, quality masks, and spatial targets.}
  \Description{Five observations, four transport residuals, quality masks,
  and spatial contribution targets for one sample.}
  \label{fig:sample-construction}
\end{figure}

\subsection{Dataset Stratification}

To prevent dataset leakage, we group overlapping AOIs into clusters
and strictly assign every AOI in a common cluster to the same dataset split. A mixed-integer optimization
assigns clusters to the training, validation, and test sets by minimizing deviations from 70\%, 15\%,
and 15\% split targets for both AOI and sample counts. After fixing the training clusters, we divide the remaining clusters
between validation and test while also balancing their class counts.

Steady class samples are more common than emission changes, and increase samples are more
common than decrease samples. We therefore downsample each training class to the size
of the smallest class. For validation and test, we retain every increase and
decrease sequence and sample steady sequences to match the larger change class.
This preserves the imbalance between increases and decreases in the held-out
sets while preventing steady behavior from dominating evaluation.

\begin{figure}[t]
  \centering
  \IfFileExists{figures/dataset-stratification.pdf}{%
      \includegraphics[width=\columnwidth]{figures/dataset-stratification.pdf}%
  }{%
      \fbox{\parbox[c][1.4in][c]{0.9\columnwidth}{\centering
      \textit{Dataset split and class-count summary}}}
  }
  \caption{AOI and sample counts by split, with class composition before and
  after sampling.}
  \Description{Training, validation, and test sizes and their decrease,
  steady, and increase class counts before and after sampling.}
  \label{fig:dataset-stratification}
\end{figure}

\bibliographystyle{ACM-Reference-Format}
\bibliography{references}

\end{document}
